%% GENERATED by `pipeline run results` - do not edit by hand.
\begin{table}[htbp]
\centering
\caption{What the validity indices actually track.}\label{tab:index-mechanism}
\footnotesize
\begin{tabular}{lll}
\toprule
relation & population & $\rho$ \\
\midrule
Silhouette vs.\ share of cases discarded as noise & all 14 & -0.308 \\
Share discarded vs.\ accuracy on the partition's own set & all 14 & -0.204 \\
Silhouette vs.\ median cluster size & full-cov.\ 14 & 0.427 \\
Silhouette vs.\ number of clusters & full-cov.\ 14 & -0.325 \\
Number of clusters vs.\ coverage & full-cov.\ 14 & -0.991 \\
Silhouette vs.\ coverage & full-cov.\ 14 & 0.341 \\
\textbf{Silhouette vs.\ common-set accuracy} & full-cov.\ 14 & \textbf{-0.392} \\
\textbf{DBCV vs.\ common-set accuracy} & full-cov.\ 14 & \textbf{-0.057} \\
Silhouette vs.\ accuracy, small clusters only (mean top-1 0.623) & 233 clusters, 14 values & -0.151 \\
Silhouette vs.\ accuracy, medium clusters only (mean top-1 0.655) & 267 clusters, 14 values & -0.009 \\
Silhouette vs.\ accuracy, large clusters only (mean top-1 0.781) & 89 clusters, 14 values & 0.074 \\
\bottomrule
\end{tabular}
\par\vspace{2pt}\footnotesize\raggedright Granularity is the relationship that survives scrutiny at this sample size. The coverage rows are largely a restatement of it (partial correlation given median cluster size is small), and the difference between the coverage and accuracy correlations is not resolvable here -- its bootstrap interval over configurations includes zero. The band rows carry only as many distinct silhouette values as there are partitions, not as many as there are clusters, so they are far weaker evidence than their cluster counts suggest.
\end{table}
